\documentclass[11pt, twocolumn]{article}
\usepackage[margin=0.65in, columnsep=0.25in, bottom=0.45in]{geometry}
\usepackage{newtxtext}
\usepackage{titlesec}
\usepackage{enumitem}
\usepackage{hyperref}
\usepackage{xcolor}

% Tight compacting
\setlist{nosep, leftmargin=*} 
\linespread{0.93} 
\setlength{\parskip}{2.5pt} 
\renewcommand{\footnotesize}{\tiny}

% Tight section formatting with visual distinction
\titleformat{\section}{\large\bfseries\scshape}{\thesection}{1em}{}[\vspace{1pt}\titlerule]
\titlespacing*{\section}{0pt}{6pt}{3pt}

\hypersetup{
    colorlinks=true,
    linkcolor=black,
    urlcolor=blue,
    citecolor=black,
    breaklinks=true,
}

\pagestyle{plain} 

% Compact Bibliography
\makeatletter
\renewenvironment{thebibliography}[1]
     {\section*{References}%
      \tiny
      \setlength{\topsep}{0pt}
      \list{\@biblabel{\@arabic\c@enumiv}}%
           {\settowidth\labelwidth{\@biblabel{#1}}%
            \leftmargin\labelwidth
            \advance\leftmargin\labelsep
            \itemsep0pt\parsep0pt\parskip0pt
            \usecounter{enumiv}%
            \let\p@enumiv\@empty
            \renewcommand\theenumiv{\@arabic\c@enumiv}}%
      \sloppy
      \clubpenalty4000
      \@clubpenalty \clubpenalty
      \widowpenalty4000%
      \sfcode`\.\@m}
     {\def\@noitemerr
       {\@latex@warning{Empty `thebibliography' environment}}%
      \endlist}
\makeatother

\begin{document}

\twocolumn[
  \begin{@twocolumnfalse}
    \noindent
    \textbf{\LARGE AI Reading Exercise: Class 11} \hfill \textbf{Daniel Hardesty Lewis} \\
    \noindent
    \textit{AI and Public Safety} \hfill April 6, 2026 \\
    \textbf{PLAN A6613: AI and the Future of Cities} \hfill Spring 2026
    \vspace{3pt}
    \hrule
    \vspace{6pt}
  \end{@twocolumnfalse}
]

\section*{Tool \& Process}
I utilized \textbf{Gemini 3.1 Pro (High)} configured through the Antigravity IDE, deploying the cutting-edge \textit{V4 Framework}.\footnote{Artifacts including the pre-audit source matrix and the bifurcated empirical/semantic audit table are archived at \url{https://github.com/dhardestylewis/plan-a6613-ai-reading-class-3/tree/main/week11}.} To prevent narrative drift and source echo-chambers, the AI performed a pre-audit on three triangulated stakeholders, bifurcated its extraction into empirical strictness and semantic tracing, and crucially subjected the resulting consensus to an \textit{Adversarial External Audit} via NIST benchmarking.

\section*{Key Findings}
\noindent Modern law enforcement agencies are aggressively deploying facial recognition technology (FRT) as a rapid biometric matching mechanism to clear investigative backlogs. However, operational realities demonstrate that deploying these proprietary, algorithmically opaque systems systematically collapses standard oversight mechanisms.

\textbf{\textit{1. Vendor claims of efficiency mask operational fragility.}}
Industry vendors heavily market FRT as a neutral biometric triage solution. They highlight objective mathematism, emphasizing how software measures facial nodal points to rapidly extract supposedly flawless biometric signatures (DataWorks Plus, 2024). Crucially, vendors defend the technology against downward liability by framing it as a strictly subordinate tool. They argue the software merely provisions a ranked list of potential identities, theoretically preserving human autonomy via a ``human-in-the-loop'' safeguard designed to adjudicate the final match (DataWorks Plus, 2024).

\textbf{\textit{2. The ``Human-in-the-Loop'' safeguard actively fails.}}
In application, this theoretical safeguard is wholly insufficient. While explicit civil liberties advocacy claims that FRT platforms inherently encode severe systemic racial disparities (ACLU, 2024), an adversarial test against federal standards provides crucial nuance. The National Institute of Standards and Technology (NIST) reports that top-tier algorithms have practically "undetectable" demographic differentials in ideal conditions (NIST, 2025). The danger, therefore, lies entirely in police operational deployment. When municipalities feed low-quality CCTV imagery into the system, false-positive error rates spike severely. 

When investigators review these algorithmically ranked outputs, they are subject to automation bias---an inherent psychological deference to the machine's judgment over independent evidence. This environmental fragility compounded by automation bias precipitated the catastrophic wrongful arrest of Robert Williams in Detroit in 2020. Because the machine produced a "match," human oversight evaporated. Williams, a Black man, was wrongfully arrested at his home and forced to remain ``locked in a crowded, filthy cell overnight'' (ACLU, 2024). This failure of human oversight ultimately cost the city a massive \$300,000 public settlement (ACLU, 2024).

\textbf{\textit{3. Binding policy must explicitly degrade algorithmic authority.}}
In response to the settlement and public outrage, the Detroit Police Department's (DPD) operational framework was radically restricted. To bypass automation bias, formal institutional policy now mandates that an ``FRT 'match' [cannot be] the sole basis for an arrest'' and explicitly dictates that ``under no circumstances does a facial recognition search result inherently establish probable cause'' (DPD, 2024). 

Furthermore, the revised policy actively dismantles the ability of algorithms to contaminate external evidence. Officers are ``strictly prohibited from... Utilizing an FRT match image in a photographic array or lineup'' to prevent a downstream cognitive cascade among civilian eyewitnesses (DPD, 2024). Finally, DPD personnel utilizing the database ``must complete specialized training on algorithmic bias, particularly the elevated error rates associated with low-quality imagery'' (DPD, 2024), formally recognizing the gap between vendor lab performance and real-world deployment.

\section*{Verification and Reflection}
The V4 Adversarial Audit prevented a critical analytical error. Unchecked, the AI previously adopted the ACLU’s narrative that the technology is uniformly and mathematically racist. The NIST external audit proved this technically false for modern top-tier models, redirecting the synthesis to correctly identify the hazard: the danger is not exclusively in the mathematics, but in the intersection of low-quality operational inputs (CCTV) and catastrophic human automation bias. This framework proves that to achieve graduate-level analysis, AI systems must be structurally forced to challenge the boundaries of the primary texts they are fed.

\begin{thebibliography}{9}
\bibitem{aclu2024} American Civil Liberties Union of Michigan. (2024). \textit{Historic Settlement Reached in Defective Facial Recognition Police Wrongful Arrest Case}.
\bibitem{dpd2024} Detroit Police Department. (2024). \textit{Manual: Facial Recognition Technology (Directive 304.1)}.
\bibitem{nist2025} National Institute of Standards and Technology. (2025). \textit{Face Recognition Technology Evaluation (FRTE) Demographic Differentials}.
\bibitem{dataworks2024} DataWorks Plus. (2024). \textit{FACE Plus Facial Recognition Overview}. \url{https://www.dataworksplus.com/faceplus.html}
\end{thebibliography}

\end{document}
